Clinical prediction models are used to estimate an individual's risk of an event, and the usefulness of that estimate depends on both how well it ranks patients (discrimination) and whether a predicted risk of, say, 5\% corresponds to about 5\% of such patients having the event (calibration)~\cite{guo2017calibration}. Rare outcomes make the training data imbalanced, and it is common practice to ``correct'' this by reweighting classes, oversampling the minority class, synthesising minority examples with SMOTE~\cite{chawla2011smote}, or moving the decision threshold. Such corrections are widely available in software~\cite{lemaitre2016imbalanced,pedregosa2012scikit}, but they alter the prevalence seen by the learner, so the output of a corrected model is not a risk estimate for the target population unless it is mapped back~\cite{pozzolo2015calibrating}.

Recent simulation and empirical studies report that corrections do not improve and can harm clinical prediction models. Carriero et al.\ studied the effect of corrections on calibration for a variety of machine learning algorithms in Monte Carlo simulations that vary sample size, predictors and event fraction, and found over-estimation of risk that re-calibration could not always remove~\cite{carriero2024harms}; Andersen et al.\ evaluated SMOTE, random undersampling and ROS at 1:1 on ten real clinical datasets with several model families, with discrimination and calibration~\cite{andersen2026tipping}; and Sirikul et al.\ studied penalised logistic regression on GUSTO-I-derived data and report that corrections failed to enhance discrimination, calibration and prediction stability~\cite{sirikul2026class}. What we add is narrow: a check on a dataset bundled with scikit-learn that anyone can re-run in minutes, with threshold moving as the arm that cannot change probabilities and the analytic prior-shift offset as a candidate repair.

We study the seed question ``which corrections help discrimination and which only distort predicted risk?'' with four hypotheses registered before the main runs (H1--H4, Section~\ref{sec:setup}). \textbf{Contributions:} (i) a controlled comparison of five strategies, two learners and four prevalences on discrimination, calibration and operating point, with 10 seeds per cell; (ii) an ablation of the prior-shift offset and a sweep of the SMOTE target ratio; (iii) an account of where the comparison is uninformative. We claim no new method.
